\begin{proposition}[The exact carry polygon does not detect Ap\'ery zeros]
\label{prop:apery-mellin-carry-polygon}
Let $p$ be an odd prime, let $0\leq j\leq p-1$, and define Morita's
$p$-adic gamma function at positive integers by
\[
 \Gamma_p(n)=(-1)^n
 \prod_{\substack{1\leq r<n\\p\nmid r}}r.
\]
For $0\leq k\leq j$, put
\[
 \epsilon_{j,k}=\mathbf 1_{\{j+k\geq p\}},
 \qquad
 U_{p,j,k}=
 \frac{\Gamma_p(j+k+1)}
 {\Gamma_p(k+1)^2\Gamma_p(j-k+1)}\in\mathbb Z_p^\times.
\]
Then one has the exact identity in $\mathbb Z_p$
\begin{equation}
 \boxed{
 b_j=\sum_{k=0}^j p^{2\epsilon_{j,k}}U_{p,j,k}^2.}
 \label{eq:apery-Morita-exact}
\end{equation}
Consequently, if
\[
 u_j=\min(j,p-1-j)+1,
 \qquad
 h_j=j+1-u_j=\max(0,2j+1-p),
\]
then the sorted valuations of the $j+1$ summands in
\eqref{eq:apery-Morita-exact} are
\begin{equation}
 \underbrace{0,\ldots,0}_{u_j\ \mathrm{times}},
 \quad
 \underbrace{2,\ldots,2}_{h_j\ \mathrm{times}}.
 \label{eq:apery-carry-slopes}
\end{equation}
In particular, every interior index $1\leq j\leq p-2$ has at least two
minimal-valuation summands.

Reduction of~\eqref{eq:apery-Morita-exact} gives the exact unit-residue
criterion
\begin{equation}
 p\mid b_j
 \quad\Longleftrightarrow\quad
 \Theta_{p,j}:=
 \sum_{k=0}^{\min(j,p-1-j)}\overline{U}_{p,j,k}^{\,2}=0
 \quad\text{in }\mathbb F_p.
 \label{eq:apery-unit-residue-criterion}
\end{equation}
The valuation data in~\eqref{eq:apery-carry-slopes} do not determine this
event.  Indeed, at $j=5$ the six summands have valuation vector
$(0,0,0,0,0,0)$ both for $p=11$ and for $p=13$, whereas their respective
residue vectors are
\[
 (1,9,1,1,9,1)\in\mathbb F_{11}^6,
 \qquad
 (1,3,4,1,10,12)\in\mathbb F_{13}^6.
\]
The first vector sums to zero and the second to~$5$.  Thus
$11\mid b_5$ but $13\nmid b_5$ despite the identical complete carry
polygon.
\end{proposition}

\begin{proof}
For $0\leq N\leq2p-2$, there is at most one multiple of~$p$ in $N!$,
and the definition of $\Gamma_p$ gives
\[
 N!=(-1)^{N+1}p^{\lfloor N/p\rfloor}\Gamma_p(N+1).
\]
The elementary cancellation
\[
 \binom jk\binom{j+k}{k}
 =\frac{(j+k)!}{k!^2(j-k)!}
\]
and the preceding factorial identity therefore give
\[
 \binom jk\binom{j+k}{k}
 =p^{\epsilon_{j,k}}U_{p,j,k};
\]
all signs cancel because the numerator and denominator sign exponents
differ by~$2$.  Squaring and summing over~$k$ proves
\eqref{eq:apery-Morita-exact}.

The no-carry range is $0\leq k\leq\min(j,p-1-j)$, which has $u_j$
elements.  Every other summand has exactly one factorial carry and hence
valuation~$2$.  This proves~\eqref{eq:apery-carry-slopes}.  Since
$u_j\geq2$ for $1\leq j\leq p-2$, the minimal valuation is never unique
in the interior.  Reducing the exact identity modulo~$p$ proves
\eqref{eq:apery-unit-residue-criterion}.  Finally, direct reduction of the
six integer summands
\[
 \binom5k^2\binom{5+k}{k}^2\qquad(0\leq k\leq5)
\]
gives the two displayed vectors and completes the counterexample.
\end{proof}

\begin{remark}[Scope of the obstruction]
\label{rem:apery-carry-polygon-scope}
The polygon in~\eqref{eq:apery-carry-slopes} is the exact termwise
Gross--Koblitz--Stickelberger carry polygon of the finite hypergeometric
sum.  It is not being identified with the Newton polygon of Frobenius on a
Kummer-twisted rigid cohomology group.  The proposition proves that this
termwise polygon, even with every summand labelled, cannot decide the
coefficient-zero event.  A successful crystalline argument must control
the reduction of the normalized slope-zero trace
$\Theta_{p,j}$, rather than only its slopes, rank, determinant, or
conductor.
\end{remark}

\begin{remark}[Finite audit]
\label{rem:apery-carry-polygon-audit}
An exact audit of all primes $p\leq200$, independently comparing the
division-free recurrence with the defining binomial sum, finds
\[
 \max_{p\leq200}Z(p)=4,
\]
attained at $p=181$ and~$193$.  No nonwrapping interval of width
$\sqrt p$ contains more than two zeros in this range.  This is a finite
certificate only; the carry polygon gives no asymptotic improvement over
the available zero-count bounds.
\end{remark}
